
\medskip
\phantomsection\label{cor:sga3-VIB-5.7.1}%
\noindent\textbf{Corollary 5.7.1.---}
\emph{Let \(S\) and \(G\) be as in
\hyperref[thm:sga3-VIB-5.7]{5.7}, and let \(\sigma,\tau\) be two
\(S\)-sections of \(G^0\), that is,
\(\sigma,\tau\in G^0(S)\). Then the coincidence subscheme
\(S(\sigma,\tau)\subset S\) of \(\sigma\) and \(\tau\)
\textup{(that is, the inverse image of the diagonal of
\(G\times_S G\) under the morphism \((\sigma,\tau)\))} is closed.}

Indeed, for every \(s\in S\), let \(U\) be an affine open subset
of \(G\) containing \(\epsilon(s)\). Then
\[
  V=\epsilon^{-1}(U)
\]
is an open neighborhood of \(s\) in \(S\), and, since \(G^0U\)
is separated, \(S(\sigma,\tau)\cap V\) is closed in \(V\).

\medskip
\phantomsection\label{rem:sga3-VIB-5.7.2}%
\noindent\textbf{Remark 5.7.2.---}
Gabber points out that one can show the following: if \(S\) is
henselian local, with closed point \(s\), then the intersection
of the open subsets \(G^0U\), where \(U\) ranges over the affine
open neighborhoods of \(\epsilon(s)\), is an open group
subscheme of \(G\), separated over \(S\).

\section*{5.8. Complements}
\phantomsection\label{sec:sga3-VIB-5.8}%
\addcontentsline{toc}{section}{\SGARefLink{sec:sga3-VIB-5.8}{5.8}. Complements}

\noindent\footnote{\label{note:sga3-VIB-47}Editor's note: we have
added this paragraph of complements and counterexamples, all
communicated by O.\ Gabber.}%
We begin by recalling
\hyperref[prop:sga3-VIA-2.6.6]{Proposition VI\(_A\), 2.6.6}:

\medskip
\phantomsection\label{prop:sga3-VIB-5.8.1}%
\noindent\textbf{Proposition 5.8.1.---}
\emph{Let \(k\) be a field, let \(G\) be a \(k\)-group locally of
finite type acting on a \(k\)-scheme \(X\) in such a way that the
morphism}
\[
  G\times X\longrightarrow X\times X,
  \qquad (g,x)\longmapsto(gx,x),
\]
\emph{is surjective. Suppose that \(X\) is quasi-separated. Then
the connected components of \(X\) are irreducible.}

\medskip
\phantomsection\label{cor:sga3-VIB-5.8.2}%
\noindent\textbf{Corollary 5.8.2.---}
\emph{Let \(S,G,X\) be as in the preliminary hypotheses of
\hyperref[thm:sga3-VIB-5.3]{5.3}. Suppose moreover that every
fiber \(X_s\) is quasi-separated and connected. Then \(X\) is
separated and quasi-compact over \(S\).}

Indeed, by the preceding proposition, every fiber \(X_s\) is
irreducible, and the rest of the proof is identical to that of
\hyperref[cor:sga3-VIB-5.4]{5.4}.
